% Lời thoại tiếng Việt cho từng slide, viết theo lối KỂ CHUYỆN: câu ngắn, ví dụ dễ hình dung, số tròn,
% không đọc tên lớp/hàm. Tên kỹ thuật chỉ nằm ở phần "Nhớ thêm" để tra khi thầy hỏi.
% \slide[must]{số}{ý chính một câu}{lời nói}{nhớ thêm (các \item, hoặc để trống)}
% \cue{...} là thao tác, không đọc lên; \luocbo{...} là đoạn nói thêm nếu còn giờ (bỏ đi vẫn liền mạch).
% guide/gen_facts.py kiểm tra file này có đủ lời nói cho đúng các slide của file PowerPoint
% và tính thời gian nói từ độ dài lời nói.

\subsection{Mở đầu (slide 1--4)}

% T1: Title slide
\slide[must]{1}{Họa. là một cửa hàng tranh biết nghe, biết nhìn; cả bài nói là hành trình đưa nó từ ý tưởng đến sản phẩm chạy thật.}{
Kính chào thầy cô và các bạn. Nhóm em xin kể về một cửa hàng tranh có tên là Họa.

Các bạn thử tưởng tượng: bạn thích một bức tranh nhưng không nhớ tên, cũng không biết họa sĩ là ai. Ở Họa., bạn có thể gõ vài chữ để tả nó. Hoặc bạn nói thẳng vào micro. Hoặc bạn đưa luôn một tấm ảnh. Thậm chí bạn chụp căn phòng của mình, và cửa hàng sẽ chọn tranh hợp với phòng, rồi treo thử lên tường cho bạn xem.

Đề tài của nhóm là từ yêu cầu đến code chạy được. Nghĩa là nhóm em đi từ câu hỏi cửa hàng này cần làm gì, cho đến một sản phẩm chạy thật, và các bạn sẽ thấy mọi bước đều khớp với nhau. \cue{chỉ vào bức tranh nền} Bức tranh nền là The Great Wave của họa sĩ Hokusai, ảnh lấy từ Wikimedia Commons. Xin phép được giới thiệu nhóm.
}{
\item ``Họa'' nghĩa là vẽ, như trong ``họa sĩ''. Nếu thầy hỏi vì sao đặt tên vậy, đó là câu trả lời.
\item Nguồn và giấy phép ảnh tranh ghi trong \file{demo/data/ATTRIBUTION.md}, phòng khi thầy hỏi về bản quyền.
\item Nói chậm, nhìn người nghe. Câu hỏi ``bạn thử tưởng tượng'' là để người nghe vào chuyện ngay.
}

% T2: Group 10 — the people behind Họa.
\slide[must]{2}{Giới thiệu nhóm 10 và ai nói phần nào.}{
Nhóm em là nhóm 10, gồm bốn thành viên: Bế Quốc Khánh, Vũ Hùng Anh, Nguyễn Gia Đức Trung và Phạm Tuấn Anh. \cue{chỉ vào khung ảnh của từng bạn khi nhắc tên}

\todo{Mỗi bạn tự nói một câu về mình, ví dụ: em là \ldots, em phụ trách phần \ldots}

Hôm nay nhóm chia thế này: bạn \todo{tên Người 1} kể phần mở đầu, Chương một và Chương hai; bạn \todo{tên Người 2} kể Chương ba; bạn \todo{tên Người 3} kể Chương bốn; và bạn \todo{tên Người 4} kể Chương năm, làm phần demo và kết luận. Xin phép bắt đầu.
}{
\item Điền tên thật vào các chỗ đỏ (xem phân công ở mục~\ref{sec:tongquan}). Mỗi bạn chỉ một câu, không kể lể.
}

% T3: One system, five perspectives
\slide{3}{Bài nói có năm chặng, mỗi chặng là một chương báo cáo, nối với nhau thành một sợi chỉ.}{
Bài nói của nhóm đi qua năm chặng, mỗi chặng là một chương của báo cáo. Chặng một: cửa hàng cần làm gì, cho ai. Chặng hai: vẽ ra cách nó hoạt động. Chặng ba: phần AI, phần nhóm em thấy thú vị nhất. Chặng bốn: biến bản vẽ thành thiết kế phần mềm. Chặng năm: code thật và một buổi demo.

\cue{chỉ vào dòng chữ ở dưới cùng} Dòng chữ ở dưới là sợi chỉ mà thầy đã dạy: đi từ vấn đề, đến yêu cầu, đến thiết kế, rồi đến code. Mỗi chặng là một mắt xích, và nhóm em sẽ cho thấy các mắt xích này nối với nhau, không bị đứt chỗ nào.
}{
\item Đọc lướt, không dừng quá 20 giây.
}

% T4: Họa. — a painting shop that listens and sees
\slide[must]{4}{Họa. bán tranh thật, có năm điều AI, và các con số trên slide được đọc thẳng từ một file chung nên luôn khớp.}{
Họa. bán bản in của \factPaintings{} bức tranh nổi tiếng thuộc phạm vi công cộng, cùng \factFrames{} loại khung, giá bằng đô la. Ngoài những việc quen thuộc như bỏ vào giỏ, thanh toán, theo dõi và hủy đơn, Họa. có \key{năm điều AI}: hiểu câu khách gõ, nghe giọng nói, nhìn ảnh một bức tranh, nhìn ảnh căn phòng để chọn tranh, và một trợ lý trả lời về đơn hàng.

\luocbo{Còn người quản lý cửa hàng thì đổi được giá, số hàng trong kho và cập nhật đơn.}

\cue{chỉ vào dòng số ở cuối slide} Có một chi tiết nhỏ mà nhóm em muốn các bạn để ý: những con số ở cuối slide, \factReq{} yêu cầu, \factUC{} việc người dùng làm được, \factJavaClasses{} lớp code, \factTables{} bảng dữ liệu, không phải nhóm em gõ tay. Chúng được đọc thẳng từ một file chung của dự án, nên luôn khớp với báo cáo, với sơ đồ và với code. Đó cũng là điều nhóm em muốn gửi gắm trong cả bài: mọi thứ phải khớp nhau.
}{
\item Công nghệ nếu thầy hỏi: Java và SQLite cho phần chính, HTML, CSS, JavaScript cho giao diện, Gemma 4 và Whisper cho AI.
\item Ảnh chụp trang chủ bên phải để cho thấy đây là cửa hàng thật, không chỉ là sơ đồ.
}

\subsection{Chương 1 -- Xác định yêu cầu (slide 5--10)}

% T5: Chapter 1 - Requirement Determination
\slide{5}{Chương 1: cửa hàng phải làm gì, và làm cho ai.}{
Chương một trả lời hai câu hỏi đơn giản: cửa hàng phải làm gì, và làm cho ai? Trước khi nghĩ tới công nghệ, nhóm em hiểu cửa hàng đã.
}{}

% T6: Understand the business before drawing diagrams
\slide{6}{Một hệ thống thông tin là con người, cách làm việc, dữ liệu và công nghệ; và ta cần vẽ ranh giới quanh nó.}{
Thầy dạy rằng một hệ thống thông tin không chỉ là phần mềm. Nó gồm bốn thứ. Con người: ở Họa. là khách, người quản lý, nhân viên kho và giao hàng. Cách làm việc: tìm tranh, bỏ vào giỏ, thanh toán, giao hàng, hủy khi cần. Dữ liệu: tranh, họa sĩ, giỏ, đơn, thanh toán. Và công nghệ, để làm cho các thứ kia chạy.

\cue{chỉ vào bảng bên trái} Bảng bên trái liệt kê từng người liên quan: họ muốn gì và cần biết gì. \cue{chỉ vào sơ đồ bên phải} Sơ đồ bên phải vẽ một đường ranh giới quanh cửa hàng: bên trong là Họa., bên ngoài là bốn bên mà nó làm việc cùng, gồm khách, người quản lý, dịch vụ thanh toán và Gemma.

Có một chi tiết hay: Gemma nằm ngoài ranh giới vì chạy trên một máy chủ khác, còn chương trình nghe giọng nói Whisper nằm trong, vì chạy ngay trên máy của cửa hàng. Nhóm em chọn vậy để giọng nói của khách không phải gửi đi đâu cả.
}{
\item Tên nếu thầy hỏi: context diagram, stakeholder, NFR08 (giọng nói chỉ xử lý trên máy).
}

% T7: From business problem to system requirement
\slide[must]{7}{Từ một vấn đề rất đời thường, đi từng bước tới yêu cầu: vấn đề, nguyên nhân, nhu cầu, yêu cầu, việc cần làm. Nhu cầu không phải là giải pháp.}{
\cue{hỏi người nghe} Có ai từng thích một bức tranh mà không biết tên nó không? Đó chính là vấn đề của khách ở Họa.

Vì sao khó tìm? Vì công cụ tìm kiếm thông thường chỉ khớp đúng từ khóa, còn hình ảnh thì rất khó tả bằng lời. Từ đó nhóm hiểu khách thật sự cần gì: tìm tranh bằng một câu nói tự nhiên, bằng giọng nói, hoặc bằng một tấm ảnh. Chỉ khi hiểu như vậy, nhóm mới viết thành yêu cầu cụ thể: hệ thống phải trả về những bức tranh giống nhất với ảnh mà khách đưa vào. Và đó trở thành một việc người dùng làm được, tên là Tìm bằng ảnh.

\cue{chỉ vào dòng chữ đỏ cuối slide} Có một điều thầy dặn rất kỹ mà nhóm em xin nhắc lại: nhu cầu không phải là giải pháp. Nói ``phải dùng Gemma'' là nói cách làm. Yêu cầu chỉ nói hệ thống phải làm được gì. Dùng công nghệ nào là chuyện của phần thiết kế, về sau.
}{
\item Câu thầy hay hỏi: ``tại sao dùng Gemma không phải là yêu cầu?''. Trả lời bằng đúng câu ``nhu cầu không phải giải pháp''.
\item Tên nếu thầy hỏi: yêu cầu FR03, use case UC04 Search by Image.
}

% T8: 14 functional requirements, prioritised with MoSCoW
\slide{8}{Mười bốn việc cửa hàng phải làm, xếp theo mức quan trọng.}{
Nhóm em liệt kê \factFR{} việc cửa hàng phải làm, rồi xếp theo mức quan trọng, đúng cách thầy dạy: bắt buộc, nên có, có thì tốt, và lần này chưa làm. Có \factMust{} việc bắt buộc, đó là lõi của cửa hàng: tìm tranh, xem chi tiết, giỏ hàng, thanh toán, xem đơn, tra cứu đơn và phần quản lý. Có \factShould{} việc ``nên có'': tìm bằng giọng nói, tìm bằng ảnh, trợ lý và hủy đơn. Có \factCould{} việc ``có thì tốt'': tìm kết hợp chữ với ảnh, và thử tranh vào phòng.

Cuối cùng, thanh toán thật là việc lần này chưa làm: phần thanh toán chỉ được mô phỏng. Nhóm em nói rõ điều này để không ai hiểu nhầm.
}{
\item MoSCoW = Must, Should, Could, Won't.
}

% T9: Non-functional requirements must be measurable
\slide{9}{Ngoài việc phải làm gì, còn việc phải làm tốt đến đâu: mỗi yêu cầu chất lượng đều có một điều kiện đo được.}{
Ngoài chuyện cửa hàng phải làm gì, còn chuyện nó phải làm tốt đến đâu. Thầy cảnh báo rằng ``hệ thống phải nhanh'' thì không đo được. Vì vậy mỗi yêu cầu chất lượng của Họa. đều đi kèm một con số hoặc một điều kiện kiểm tra được.

Ví dụ: tìm bằng chữ phải trả lời trong hai giây. Ảnh tải lên chỉ nhận ba định dạng và nhẹ hơn năm megabyte. Âm thanh chỉ được xử lý ngay trên máy này. Không bao giờ lưu số thẻ đầy đủ, chỉ lưu bốn số cuối. Khách chỉ xem được đơn của chính mình. Và nếu Gemma hỏng thì cửa hàng vẫn tìm được tranh. Tất cả có \factNFR{} yêu cầu như vậy, và nhóm em đã kiểm chứng chúng, các bạn sẽ thấy ở phần đánh giá và phần demo.
}{
\item Nếu thiếu giờ, chỉ nói hai ý: ``Gemma hỏng vẫn tìm được'' và ``khách chỉ xem được đơn của mình''. Cả hai sẽ được demo.
\item Tên nếu thầy hỏi: NFR01 bảo mật, NFR02 dự phòng, NFR03 ảnh, NFR05 thẻ, NFR07 hiệu năng, NFR08 giọng nói.
}

% T10: A use case is a goal, not a button
\slide[must]{10}{Bốn bên dùng cửa hàng, mười bốn việc họ làm được; use case là mục tiêu, không phải nút bấm.}{
Ai dùng cửa hàng? Họa. có bốn bên. Khách hàng và người quản lý là hai bên dùng trực tiếp. Gemma và dịch vụ thanh toán là hai bên mà hệ thống gọi nhờ. Mỗi việc người dùng muốn làm được gọi là một use case, và Họa. có \factUC{} việc như vậy.

Thầy nhấn mạnh hai điều. Một: người có liên quan chưa chắc là người dùng hệ thống. Nhân viên kho có liên quan, nhưng họ không bấm gì trong hệ thống nên không tính. Hai: use case là một mục tiêu, không phải một nút bấm. Ta viết ``Tìm tranh'', không viết ``bấm nút tìm''. \cue{chỉ vào sơ đồ} Trên sơ đồ, tìm bằng chữ, bằng giọng nói, bằng ảnh, và cả thử tranh vào phòng, đều là các cách khác nhau để làm cùng một việc là tìm tranh.

\luocbo{Hai quan hệ nữa: việc Đặt hàng luôn kéo theo việc Quản lý giỏ, còn việc Hủy đơn là phần thêm tùy chọn của việc Xem đơn.}
}{
\item Tên nếu thầy hỏi: actor, include, extend, generalization. Search Product viết nghiêng vì là use case trừu tượng.
\item Chỉ lần lượt trên sơ đồ: mũi tên rỗng từ bốn kiểu tìm vào Search Product; include từ Checkout sang Manage Cart; extend từ Cancel Order sang View Order.
\item Dòng cuối slide là dấu vết truy vết: FR03, UC04, hàm tìm bằng ảnh, kiểm thử.
}

\subsection{Chương 2 -- Phân tích (slide 11--15)}

% T11: Chapter 2 - Analysis
\slide{11}{Chương 2: biến yêu cầu thành các bản vẽ, lấy việc đặt hàng làm ví dụ xuyên suốt.}{
Chương hai biến các yêu cầu thành bản vẽ. Có hai câu hỏi: khi một người làm một việc, hệ thống làm gì, từng bước ra sao? Và hệ thống cần quản lý những thứ gì? Nhóm em lấy việc đặt hàng làm ví dụ xuyên suốt.
}{}

% T12: Use case specification: Checkout (UC07)
\slide[must]{12}{Việc Đặt hàng được viết ra như một kịch bản, gồm cả những lúc có chuyện.}{
Đây là việc đặt hàng, viết ra như một kịch bản. Khách đã đăng nhập và có hàng trong giỏ. Năm bước chính: cửa hàng cho khách xem giỏ và tổng tiền; khách chọn địa chỉ và cách thanh toán; cửa hàng kiểm tra lại giá và số hàng còn trong kho; tạo đơn với giá đã chốt; rồi ghi nhận thanh toán, trừ kho và làm rỗng giỏ.

Thầy dặn: chỉ viết lúc mọi chuyện suôn sẻ là chưa đủ. Nên nhóm em viết cả những lúc có chuyện. Nếu một bức hết hàng, cửa hàng nói rõ bức nào và không tạo đơn. Nếu địa chỉ không phải của khách, cửa hàng yêu cầu chọn lại. Nếu thanh toán bị từ chối, khách chọn cách khác. \luocbo{Trong báo cáo còn một tình huống nữa: khách gõ một địa chỉ mới ngay lúc đặt hàng, và nhóm có làm trong phần demo.}
}{
\item Tên nếu thầy hỏi: UC07 Checkout, luồng chính, luồng thay thế.
}

% T13: From scenario to activity: decisions follow rules
\slide{13}{Từ kịch bản vẽ thành sơ đồ từng bước; mỗi ngã rẽ trong sơ đồ đến từ một luật của cửa hàng.}{
Từ kịch bản, nhóm vẽ thành một sơ đồ chạy từng bước, chia làm ba làn: khách, cửa hàng và cơ sở dữ liệu. Chỗ đáng chú ý là mỗi ngã rẽ trong sơ đồ đều đến từ một luật của cửa hàng: hết hàng thì không tạo đơn; địa chỉ phải là của khách; chỉ được hủy khi đơn chưa giao.

Bên phải là bốn câu hỏi thầy dạy để phân tích một việc: người dùng muốn gì, bình thường diễn ra thế nào, điều gì có thể sai, và hệ thống làm gì khi đó. Nhờ vậy ta luôn lần được từ một yêu cầu tới việc cần làm, tới kịch bản, rồi tới sơ đồ.
}{
\item Tên nếu thầy hỏi: activity diagram, ba swimlane.
}

% T14: Consistency check: cancel before shipment
\slide{14}{Kiểm tra các bản vẽ có mâu thuẫn nhau không bằng cách lần theo một yêu cầu: ``được hủy đơn trước khi giao''.}{
Để chắc các bản vẽ không mâu thuẫn nhau, nhóm em làm đúng bài thầy cho: lần theo yêu cầu ``khách được hủy đơn trước khi giao''. Yêu cầu này xuất hiện ở năm nơi, và cả năm đều khớp. Trong danh sách yêu cầu, có yêu cầu hủy đơn. Trong các việc người dùng làm được, có việc Hủy đơn. Trong sơ đồ từng bước, có câu hỏi ``đơn đã giao chưa?''. Trong sơ đồ trạng thái của đơn, chỉ cho hủy khi đơn mới đặt hoặc đang xử lý. Và trong thiết kế, hàm hủy đơn kiểm tra đúng luật đó.

Còn phía quản lý: khi đơn chuyển sang trạng thái đã gửi, khách hết hủy được. \cue{chỉ vào ảnh bên phải} Ảnh bên phải là chức năng này chạy thật, nhóm em sẽ làm lại trong phần demo.
}{
\item Tên nếu thầy hỏi: FR11, UC11 mở rộng UC08, máy trạng thái PLACED, PROCESSING, SHIPPED, DELIVERED, hàm cancelOrder gọi Order.cancel.
}

% T15: Domain concepts: CartItem ≠ OrderItem
\slide[must]{15}{Giỏ hàng khác đơn hàng: giỏ là danh sách đang chọn, đơn là hóa đơn đã chốt giá; đặt hàng là biến giỏ thành đơn.}{
Đây là bức tranh những thứ cửa hàng cần quản lý: khách, tranh và khung, họa sĩ, giỏ, đơn, vận chuyển, thanh toán. Nhóm có \factAnalysis{} khái niệm như vậy.

Điều thầy nhấn mạnh, và nhóm em thích nhất, là món trong giỏ và món trong đơn là hai thứ khác nhau. Giỏ hàng giống một danh sách đang chọn: giá có thể đổi, khách có thể bỏ bớt. Đơn hàng giống một hóa đơn: giá đã chốt, không đổi nữa. Cho nên đặt hàng thực chất là biến một giỏ thành một đơn. \cue{chỉ vào mũi tên đứt nét giữa giỏ và đơn} Mũi tên đứt nét này thể hiện đúng điều đó: giỏ tạo ra đơn, rồi giỏ rỗng đi.

\luocbo{Mỗi khách có tối đa một giỏ, một giỏ có nhiều món, và một đơn có ít nhất một món.}

Hết phần phân tích. Xin nhường lời cho bạn \todo{tên Người 2} kể Chương ba, phần AI.
}{
\item Tên nếu thầy hỏi: CartItem khác OrderItem, quan hệ phụ thuộc ``create''.
\item Giá trong giỏ có thể cũ; lúc đặt hàng cửa hàng đọc lại giá hiện tại rồi mới chốt.
}

\subsection{Chương 3 -- AI trong thương mại điện tử (slide 16--23)}

% T16: Chapter 3 - AI in e-commerce
\slide{16}{Năm điều AI của Họa.; nguyên tắc: AI giúp hiểu khách, còn giá, kho và quyền xem đơn do luật của cửa hàng quyết định.}{
Em xin tiếp tục với Chương ba, phần mà nhóm em thấy thú vị nhất: AI. Họa. có năm điều AI: hiểu câu khách gõ, nghe giọng nói, nhìn ảnh một bức tranh, nhìn ảnh căn phòng để chọn tranh, và một trợ lý mua sắm.

Nhóm em đặt ra một nguyên tắc cho cả năm điều này: AI giúp khách nói ra điều mình muốn, còn giá, số hàng trong kho và quyền xem đơn thì luôn do luật của cửa hàng quyết định, không phải do AI.
}{}

% T17: Different inputs, one common query
\slide[must]{17}{Dù khách tìm bằng cách nào, cửa hàng cũng làm theo một đường chung: đưa về một dạng chung, chọn ứng viên, xếp hạng.}{
Dù khách tìm bằng cách nào, bằng chữ, giọng nói, ảnh tranh hay ảnh phòng, cửa hàng đều đi theo một đường chung, đúng như bài tập của thầy: đưa mọi thứ về một dạng chung, chọn ra những bức có thể hợp, rồi xếp hạng.

\cue{chỉ lần lượt bốn hộp bên trái} Chữ thì đi thẳng vào. Giọng nói thì được đổi thành chữ trước. Ảnh một bức tranh thì được Gemma tả thành vài từ khóa. Ảnh căn phòng thì được Gemma đọc ra màu sắc và phong cách. Tất cả gặp nhau ở một phiếu yêu cầu chung, rồi cùng đi qua bước chọn và bước xếp hạng, và cuối cùng là danh sách tranh trả về cho khách.

\luocbo{Nhờ có một đường chung nên thêm một cách tìm mới dễ hơn: chỉ cần thêm một cách đổi đầu vào thành phiếu yêu cầu.}
}{
\item Tên nếu thầy hỏi, theo từng hộp: SearchUI.submitText, QueryService.textQuery; VoiceInput.upload, SpeechService.transcribe, QueryService.voiceQuery; ImageUpload.choose, ImageService.embed, QueryService.imageQuery; RoomPreview.chooseRoom, ImageService.describe, SearchService.searchRoom.
\item Phiếu yêu cầu chung là CommonQuery; bước chọn là SearchService, bước xếp hạng là RankingService.
\item Bốn nguyên tắc ở dưới slide (A06): một biểu diễn chung; chọn tách khỏi xếp hạng; giao diện không chạm cơ sở dữ liệu; bắt đầu đơn giản.
}

% T18: Voice search: the audio never leaves the laptop
\slide{18}{Tìm bằng giọng nói: khách nói, âm thanh không rời khỏi chiếc máy này; kết quả hiện thành chữ để khách sửa.}{
Giọng nói. Khách bấm micro và nói, ví dụ ``tranh phong cảnh màu xanh''. Điều nhóm em tự hào là âm thanh không rời khỏi chiếc máy này. Nhóm dùng Whisper, một chương trình nhận dạng giọng nói, chạy ngay trên máy, không gửi lên mạng. Kết quả hiện ra thành chữ để khách xem và sửa nếu nghe sai, rồi mới đi tìm như một câu gõ bình thường. Nếu Whisper không chạy, cửa hàng chỉ nhẹ nhàng mời khách gõ, không báo lỗi khó hiểu.

\cue{chỉ vào ba con số lớn} Nhóm em đã đo thử với \evSttN{} câu tiếng Anh: trung bình \evSttMean{} giây một câu, sai khoảng \evSttWer{} số từ, và \evSttExact{} trên \evSttN{} câu đúng hoàn toàn. Nhóm cố ý không làm phần đọc thành tiếng, vì không có việc nào cần đến.
}{
\item Sơ đồ bên trái có đúng bảy thành phần như bài tập A06 yêu cầu; không cần đọc từng cái.
\item Đo bằng tiếng Anh vì sản phẩm dùng tiếng Anh. Nếu thầy hỏi vì sao slide này ghi 2\% còn slide đánh giá ghi khoảng 4\%: hai cách đo khác nhau (mục~\ref{sec:hoidap}).
\item Tên nếu thầy hỏi: SpeechService, WhisperClient, NFR08.
}

% T19: Search by image: describe it, then match it
\slide[must]{19}{Tìm bằng ảnh theo ba bước như một người bán tranh: tả ảnh bằng ba nhãn, so với các bức trong cửa hàng, rồi xếp hạng. Nói thật hạn chế.}{
Tìm bằng ảnh. Cách làm rất giống một người bán tranh. Bước một: Gemma nhìn ảnh khách đưa và nói ra ba điều: phong cách, chủ đề và màu chủ đạo. \cue{chỉ vào dòng ví dụ ở dưới} Ví dụ, đưa ảnh The Great Wave, Gemma nói: kiểu tranh Nhật ukiyo-e, vẽ biển, màu xanh.

Bước hai: mỗi bức tranh trong cửa hàng cũng được gắn sẵn đúng ba nhãn như vậy. Cửa hàng chỉ việc xem bức nào giống ảnh nhất: càng nhiều nhãn trùng thì càng giống.

Bước ba: xếp hạng. Độ khớp theo chữ chiếm bốn mươi lăm phần trăm, độ khớp theo ảnh bốn mươi lăm phần trăm, và việc còn hàng mười phần trăm. Tổng cao nhất đứng đầu.

Nhóm em xin nói thật một hạn chế: ba nhãn này do nhóm tự định ra, nên cách hiểu ảnh còn đơn giản, chưa sâu như các mô hình chuyên về ảnh. Hướng phát triển là thay bằng mô hình tốt hơn. Còn nếu Gemma không có mặt, cửa hàng chỉ nhìn được màu trung bình của ảnh, và kết quả yếu đi nhiều.
}{
\item Slide này cố ý không có công thức. Nếu thầy hỏi công thức của A06, trả lời theo mục~\ref{sec:hoidap} (câu về trọng số xếp hạng).
\item Tên nếu thầy hỏi: toVector, cosineSimilarity, RankingService.rank; trọng số 0,45, 0,45, 0,10 nằm trong file cấu hình.
\item Sơ đồ bên phải chỉ để cho thấy thứ tự các bước; không cần giải thích.
}

% T20: Try it in your room (UC14 Find Painting for My Room)
\slide[must]{20}{Thử tranh vào phòng của bạn: Gemma đọc căn phòng, chọn vài bức hợp và nói vì sao hợp; cửa hàng treo thử lên tường trong chính ảnh phòng.}{
Đây là điều nhóm em thích nhất: thử tranh vào phòng của bạn. Khách đưa một tấm ảnh căn phòng. Gemma đọc căn phòng: màu sắc, phong cách, không khí, ánh sáng. Cửa hàng chọn ra mười hai bức có vẻ hợp. Rồi Gemma xem lại ảnh phòng cùng mười hai bức đó, chọn ba đến năm bức hợp nhất, và viết cho mỗi bức một câu giải thích vì sao hợp, ví dụ: tông nâu ấm của bức này hợp với sàn gỗ và tường gạch của phòng. Cuối cùng, cửa hàng treo bức tranh lên tường ngay trong ảnh, để khách thấy trước khi mua.

Nhóm em có vài chốt chặn để Gemma không nói bừa. Gemma chỉ được chọn trong mười hai bức đã đưa ra. Câu giải thích quá ngắn hoặc chung chung sẽ bị thay bằng câu dựng từ chính thông tin của bức tranh. Và ảnh phòng chỉ nằm trong trình duyệt của khách, không lưu lên máy chủ. \cue{chỉ vào ảnh chụp} Đây là kết quả thật: bức tranh đã treo trên tường, kèm câu giải thích bên dưới.
}{
\item Tên nếu thầy hỏi: UC14, FR14 mức ``có thì tốt'', trường reason của kết quả, màn hình RoomPreview.
\item Nếu Gemma không có, thứ hạng vẫn giữ và câu giải thích được dựng từ thuộc tính của tranh.
\item Nhóm sẽ làm trực tiếp ở bước 4 của demo.
}

% T21: AI behind interfaces, always with a fallback
\slide{21}{Mọi chỗ dùng AI đều đặt sau một ``cổng'' riêng và đều có phương án dự phòng.}{
Mọi chỗ dùng AI trong Họa. đều đi qua một ``cổng'' riêng, nên có thể đổi mô hình, hoặc dùng bản giả khi kiểm thử, mà không phải sửa phần còn lại. Quan trọng hơn, mỗi chỗ đều có phương án dự phòng.

\cue{chỉ vào bảng bên phải} Gemma lỗi hoặc chậm hơn mười lăm giây thì dùng luật từ khóa. Whisper không chạy thì mời khách gõ. Ảnh sai loại hoặc quá lớn thì bị từ chối từ đầu. AI chỉ giúp hiểu khách, nó không bao giờ quyết định giá hay số hàng. Và quyền xem đơn luôn do máy chủ kiểm tra. \cue{chỉ vào ảnh bên dưới} Ảnh bên dưới là tìm kiếm vẫn chạy khi tắt Gemma, và nhóm em sẽ chứng minh trong phần demo.
}{
\item Tên nếu thầy hỏi: LlmClient cài đặt bởi GemmaClient, SpeechClient cài đặt bởi WhisperClient.
\item Bước 9 của demo chứng minh slide này.
}

% T22: Curator assistant: Gemma picks the tool
\slide{22}{Trợ lý mua sắm: khách hỏi bằng lời thường, Gemma chọn công cụ của cửa hàng để trả lời, và luôn tra theo tài khoản đang đăng nhập.}{
Trợ lý mua sắm. Khách hỏi bằng lời thường, ví dụ ``đơn gần nhất của tôi đâu rồi?''. Gemma có thể ``nhờ'' các công cụ của cửa hàng: tìm tranh, tìm đơn theo mã, xem đơn mới nhất, liệt kê đơn. Nó tự chọn công cụ nào cần dùng, rồi viết câu trả lời tự nhiên.

Điểm an toàn then chốt: cửa hàng luôn tra theo tài khoản đang đăng nhập, không theo những gì khách gõ trong câu hỏi. Nên dù khách có gõ mã đơn của người khác, cũng không xem được. Khách hỏi bằng chữ hoặc bằng giọng nói đều được; trợ lý trả lời bằng chữ.
}{
\item Tên nếu thầy hỏi: tool calling, AssistantService.ask và runTool.
\item Bước 6 và 7 của demo sẽ cho thấy điều này.
}

% T23: Evaluation as required by A06
\slide[must]{23}{Đánh giá trung thực: Gemma hơn rõ ở tìm bằng ảnh, nhưng ở tìm bằng chữ thì ngang với luật đơn giản.}{
Vậy các điều AI này tốt đến đâu? Nhóm em đo hai lần trên cùng bộ câu hỏi: một lần có Gemma, một lần tắt Gemma để chỉ còn luật từ khóa. \cue{chỉ vào biểu đồ} Cột cam là có Gemma, cột xám là không có.

Kết quả khá thú vị. Với tìm bằng chữ, hai bên ngang nhau: khoảng \evTextDirectGemma{} với câu nói thẳng, và \evTextParaGemma{} với câu diễn đạt lại. Nghĩa là với cửa hàng nhỏ này, Gemma không giúp tìm chữ tốt hơn một luật đơn giản. Nhưng với ảnh thì khác hẳn: \evImageGemma{} khi có Gemma, so với chỉ \evImageLocal{} khi không có. Giọng nói đạt \evVoiceHit{}. Và trợ lý trả lời đúng \evAssistantCorrect{} trên \evAssistantN{} câu, không lần nào làm lộ dữ liệu của người khác.

Nhóm em nói thật: đây là một lần đo, với bộ câu hỏi nhỏ, và Gemma có thể cho kết quả hơi khác giữa các lần. Vì vậy xin xem đây là gợi ý, không phải chứng minh.
}{
\item Đơn vị: tỉ lệ kết quả đầu tiên đúng (Hit@1); riêng trợ lý là tỉ lệ câu trả lời đúng. Cột giọng nói chỉ có một thanh vì không có lần chạy tắt Gemma.
\item Câu thầy hay hỏi: ``kết quả có đáng tin không?'' Trả lời theo mục~\ref{sec:hoidap}. Nhóm khẳng định đúng hai điều: ảnh thì Gemma hơn rõ; chữ thì không hơn.
\item Xin nhớ: \evVoiceWer{} sai từ ở đây khác với \evSttWer{} ở slide giọng nói vì đo cả đường ống.
\item Bảng bên phải là các câu demo, cạnh các câu ví dụ trong bài tập của thầy.
\item Hết phần AI. Xin nhường lời cho bạn \todo{tên Người 3} kể Chương bốn.
}

\subsection{Chương 4 -- Thiết kế (slide 24--28)}

% T24: Chapter 4 - Design
\slide{24}{Chương 4: phân tích nói ``cái gì'', thiết kế nói ``làm thế nào''.}{
Em xin kể Chương bốn: thiết kế. Nếu phân tích là ``cửa hàng cần gì'', thì thiết kế là ``phần mềm sẽ làm ra điều đó thế nào''. Chương này có bốn việc: biến các khái niệm thành các khối code, biến chúng thành bảng dữ liệu, chia code thành bốn tầng, và sinh code từ chính bản thiết kế.
}{}

% T25: From analysis classes (A04) to design classes (A05)
\slide[must]{25}{Biến khái niệm thành khối code: mỗi khối có thông tin và có việc làm, mỗi việc có một người phụ trách rõ ràng.}{
Việc đầu tiên: biến khái niệm thành khối code. Giỏ, món trong giỏ, đơn, món trong đơn trở thành những khối vừa có thông tin vừa có việc làm. Còn việc ``đặt hàng'' trở thành một việc cụ thể, do một người điều phối đảm nhận. Thầy dạy cách chia như vậy: một khối lo nhận yêu cầu, một khối lo điều phối, các khối nghiệp vụ lo giữ luật, và các khối lưu trữ lo cất dữ liệu.

Có một chỗ nhóm em làm khác bài giảng một chút. Việc đặt hàng trong bài chỉ nhận mã giỏ. Của nhóm còn nhận thêm mã khách lấy từ đăng nhập, để kiểm tra giỏ này đúng là của khách này, cùng với địa chỉ và cách thanh toán. \cue{chỉ vào sơ đồ bên phải} Bên phải là sơ đồ các khối thiết kế.
}{
\item Tên nếu thầy hỏi: lớp phân tích A04 sang lớp thiết kế A05, hàm CheckoutService.checkout(cartId, customerId, addressId, payment), slide 22 của bài 5.
}

% T26: From data model to an SQLite database
\slide{26}{Việc thứ hai: cất dữ liệu. Cơ sở dữ liệu được tạo thẳng từ mô hình theo ba quy tắc đơn giản.}{
Việc thứ hai là cất dữ liệu. Nhóm em tạo cơ sở dữ liệu thẳng từ mô hình, không viết tay, gồm \factTables{} bảng. Có ba quy tắc chuyển đổi chính. Một: khi một thứ là một loại của thứ khác, ví dụ khách VIP là một loại khách, thì có bảng riêng nhưng dùng chung mã với bảng cha. Hai: khi một bên có nhiều bên kia, ví dụ một khách có nhiều địa chỉ, thì bảng bên ``nhiều'' giữ mã của bên kia. Ba: khi một thứ là phần không tách rời của thứ khác, như các món trong một đơn, thì xóa đơn sẽ xóa luôn các món.

\cue{chỉ vào ảnh bên phải} Ảnh bên phải là cơ sở dữ liệu thật của sản phẩm, chính là phần ``màn hình cơ sở dữ liệu'' mà thầy yêu cầu.
}{
\item Tên nếu thầy hỏi: kế thừa thành bảng con có khóa chính cũng là khóa ngoại; một-nhiều thì khóa ngoại phía nhiều; composition dùng khóa ghép và ON DELETE CASCADE.
}

% T27: Four layers — a layer is not a package
\slide[must]{27}{Việc thứ ba: chia code thành bốn tầng, như các bộ phận của một nhà hàng, với hai luật phụ thuộc nghiêm ngặt.}{
Việc thứ ba: chia code thành bốn tầng. Cho dễ hình dung, các bạn nghĩ tới một nhà hàng. Tầng giao diện giống người phục vụ: nhận yêu cầu và mang kết quả ra. Tầng điều phối giống bếp trưởng: quyết định làm gì trước, làm gì sau. Tầng nghiệp vụ là luật của nhà hàng: món nào giá bao nhiêu, khi nào được hủy. Và tầng hạ tầng là kho, tủ lạnh, thiết bị: nơi lưu và lấy dữ liệu, nơi gọi Gemma và Whisper.

Thầy lưu ý một điều hay nhầm: tầng không phải là thư mục. Mỗi tầng được chia thành nhiều gói nhỏ, tổng cộng \factPackages{} gói. \cue{chỉ vào hai dòng ở cuối slide} Nhóm em giữ nghiêm hai luật. Giao diện không bao giờ chạm thẳng vào cơ sở dữ liệu. Và phần điều phối chỉ làm việc với ``cổng'' của kho và của AI, nên đổi thứ bên dưới thì không phải sửa thứ bên trên.
}{
\item Tên bốn tầng: presentation, application, domain, infrastructure. Các gói nằm trong ô bên phải slide.
\item Ví dụ nhà hàng chỉ để dễ hình dung; định nghĩa chính xác là: tầng là trách nhiệm kiến trúc, gói là cách tổ chức code.
}

% T28: Checkout across the layers (Lecture 05)
\slide[must]{28}{Khi khách bấm Đặt hàng, bốn tầng làm việc cùng nhau; luật nghiệp vụ nằm ở khối nghiệp vụ, không dồn hết vào người điều phối.}{
Hãy xem bốn tầng làm việc cùng nhau khi khách bấm ``Đặt hàng''. Giao diện chuyển yêu cầu cho người điều phối. Người điều phối làm lần lượt: lấy giỏ ra xem; lấy thông tin khách; kiểm tra lại giá và số hàng của từng món trong kho; tạo đơn với giá chốt; tính tổng, giảm giá nếu là khách VIP; lưu đơn; trừ kho; làm rỗng giỏ; rồi báo đặt hàng thành công. \cue{đi dọc sơ đồ từ trái sang phải}

\luocbo{Điểm hay: các luật như ``thêm một món vào giỏ'' hay ``tính tổng đơn'' nằm ở các khối nghiệp vụ, không nằm ở người điều phối. Người điều phối chỉ gọi, không ôm hết mọi việc. Thầy gọi một khối ôm hết là God class và dặn tránh.}

Thiết kế đến đây là đủ. Xin nhường lời cho bạn \todo{tên Người 4} kể Chương năm, phần code và demo.
}{
\item Tên nếu thầy hỏi: CheckoutController, CheckoutService.checkout, CartRepository, CustomerRepository, ProductRepository, OrderRepository; luật nằm ở Cart.addItem và Order.calculateTotal.
}

\subsection{Chương 5 -- Code và demo (slide 29--32)}

% T29: Chapter 5 - Code
\slide{29}{Chương 5: code không chỉ chạy được mà còn khớp với bản thiết kế từng cái tên.}{
Em xin kể Chương năm, phần code. Mục tiêu của nhóm em không chỉ là code chạy được, mà là code khớp với bản thiết kế từng cái tên: tên khối, tên việc, tên bảng, tên cột trong code phải giống hệt trong thiết kế. Phần chính viết bằng Java và SQLite, giao diện bằng HTML, CSS và JavaScript, và chương trình nghe giọng nói bằng Python.
}{}

% T30: Course code ↔ Họa. code
\slide{30}{Tên trong code giữ nguyên tên trong bài giảng của thầy, chỉ đổi theo cách đặt tên của Java.}{
Các ví dụ code mà thầy cho trong bài giảng viết bằng Python. Nhóm em giữ nguyên tên, chỉ đổi theo cách đặt tên của Java: ví dụ ``all\_products'' thành ``allProducts'', ``text\_query'' thành ``textQuery''. Nhờ vậy, ai đã đọc bài giảng thì nhìn vào code của nhóm vẫn nhận ra ngay. Cột bên phải cho biết từng khối thuộc tầng nào, khớp với bốn tầng vừa rồi.

Tổng cộng có \factJavaClasses{} khối code Java và \factJsViews{} phần giao diện, tất cả sinh ra từ cùng một file thiết kế. \luocbo{Mấy cặp tiêu biểu trong bảng: tìm bằng chữ, tìm bằng ảnh, tính độ giống, xếp hạng, tra đơn, và đặt hàng.}
}{
\item Dòng cuối slide là ánh xạ công nghệ của bài 5: bộ điều khiển là một handler HTTP của Java, kho dữ liệu là JDBC trên SQLite, cơ sở dữ liệu là file ecom.db.
\item Dòng công thức trong bảng là phần ``tính điểm xếp hạng'' của bài tập A06; không cần giải thích.
}

% T31: Representative code and how it is verified
\slide[must]{31}{Làm sao biết code thật sự khớp thiết kế? Nhóm có bốn cách tự kiểm tra.}{
Làm sao biết code thật sự khớp với thiết kế? Nhóm em có bốn cách. Một: một công cụ tự so tên giữa code và thiết kế, rồi báo tất cả đều khớp: \factCheckOps{} việc, \factCheckColumns{} cột dữ liệu và nhiều thứ khác. Hai: \factTests{} bài kiểm thử tự động, tất cả đều đạt. Ba: phần đánh giá nhóm đã trình bày. Bốn: cơ sở dữ liệu được sinh ra từ chính bản thiết kế. Nhờ vậy, báo cáo, sơ đồ và code luôn đồng bộ với nhau.

\cue{chỉ vào đoạn code bên trái} Đoạn code bên trái chỉ để minh họa: đây là phần đổi mô tả bức ảnh thành ba nhãn, rồi so xem hai bên giống nhau đến đâu.
}{
\item Tên công cụ: tools/check\_code\_model.py; kết quả nằm trong build/test\_results.txt. Con số đầy đủ: \factCheckClasses{} lớp, \factCheckFields{} trường, \factCheckOps{} hàm, \factCheckTables{} bảng, \factCheckColumns{} cột.
\item Đoạn code là VectorIndex.
}

% T32: Live demo script
\slide[must]{32}{Demo trực tiếp chín việc; mọi thứ chạy trên chiếc laptop này, trừ Gemma chạy trên một máy chủ khác.}{
Bây giờ là phần demo. Ở góc phải là sơ đồ chạy thật: mọi thứ chạy ngay trên chiếc laptop này, chỉ trừ Gemma chạy trên một máy chủ khác. Nhóm em sẽ làm chín việc: mở cửa hàng; tìm bằng chữ; tìm bằng giọng nói; tìm bằng ảnh và thử tranh vào phòng; đặt hàng; hỏi trợ lý; hủy đơn và thử xem đơn của người khác; vào trang quản lý; và cuối cùng tắt Gemma để xem cửa hàng có còn chạy không. Xin chuyển sang trình duyệt.
}{
\item Người 4 nói và chuyển sang trình duyệt; Người 3 thao tác theo mục~\ref{sec:demo}.
\item Xong demo, quay lại PowerPoint ở slide 33.
}

\subsection{Kết luận (slide 33--35)}

% T33: Traceable from requirement to code
\slide[must]{33}{Một sợi chỉ liền mạch từ yêu cầu đến code, bốn bài học, và những điều chưa làm được nói thẳng.}{
Để kết lại. Nhóm em giữ được một sợi chỉ liền mạch: từ một yêu cầu, ví dụ ``tìm tranh bằng ảnh'', qua sơ đồ, qua thiết kế, tới code, tới bảng dữ liệu, tới kiểm thử và đánh giá, không đứt chỗ nào.

Nhóm em rút ra bốn bài học. Một: hiểu vấn đề trước khi chọn công nghệ. Hai: phân tích nói ``cái gì'', thiết kế nói ``làm thế nào''. Ba: AI nên là một phần có thể thay, và luôn có phương án dự phòng. Bốn: một bản đặc tả chung giữ cho mọi thứ khớp nhau.

Còn những điều chưa làm được, nhóm em xin nói thẳng: thanh toán mới chỉ mô phỏng; bộ đánh giá còn nhỏ; gợi ý cá nhân hóa cần dữ liệu và sự đồng ý của khách; và nếu nhiều máy dùng chung thì cần cơ sở dữ liệu dạng máy chủ. \luocbo{Hai điểm kỹ thuật nữa: đặt hàng chưa nằm gọn trong một giao dịch cơ sở dữ liệu, và hủy đơn chưa hoàn lại số hàng vào kho; cả hai đều đã có cách sửa rõ ràng.}

Xin cảm ơn thầy cô và các bạn đã lắng nghe.
}{
\item Chi tiết các hạn chế và cách sửa: mục~\ref{sec:hanche}.
}

% T34: Questions & answers
\slide[must]{34}{Hỏi đáp.}{
Nhóm em xin cảm ơn thầy cô và các bạn đã lắng nghe, và sẵn sàng trả lời các câu hỏi.
}{
\item Người phụ trách chương liên quan trả lời trước (mục~\ref{sec:hoidap}). Nếu chưa chắc, mở đúng mục theo bảng ở mục~\ref{sec:filemap} thay vì đoán.
\item Hai tab trình duyệt và PDF báo cáo nên mở sẵn.
}

% T35: Folder map of tieuluan/
\slide{35}{Slide dự phòng: bản đồ thư mục, chỉ mở khi được hỏi ``cái này nằm ở đâu''.}{
Đây là slide dự phòng, dùng khi thầy hỏi một thứ nằm ở đâu trong dự án. Mỗi dòng là một thư mục: tên, nội dung bên trong, và chương nào dùng nó.
}{
\item Thông tin chi tiết hơn: mục~\ref{sec:filemap}.
}
